\documentclass[11pt]{article}
\usepackage[a4paper,margin=26mm]{geometry}
\usepackage{amsmath,amssymb,amsthm,booktabs,array}
\usepackage[hidelinks]{hyperref}
\newtheorem{theorem}{Theorem}
\newtheorem{proposition}[theorem]{Proposition}
\newtheorem{corollary}[theorem]{Corollary}
\newcommand{\Fib}{\operatorname{Fib}}
\newcommand{\NC}{\operatorname{NC}}
\newcommand{\Comp}{\operatorname{Comp}}
\newcommand{\wt}{\operatorname{wt}}
\newcommand{\hik}{\operatorname{h}}
\title{Exact support for Petrov--Scott, \S8.5,\\
Conjecture 8.22 and Problem 8.23:\\
an application and alternative proof of Flajolet (1980)}
\author{Haobo Ma\\\small ChronoAI Pte Ltd, Singapore
\and Wenlin Zhang\\\small National University of Singapore\\
\small ChronoAI Pte Ltd, Singapore}
\date{8 November 2026}
\begin{document}
\maketitle
\begin{abstract}
For Petrov and Scott's \emph{Random Fibonacci words via clone Schur
functions} (arXiv:\allowbreak2412.21126v2), \S8.5, Conjecture~8.22 and
Problem~8.23, we apply and give an alternative proof of Flajolet's
(1980, Proposition~3B) coefficient formula to obtain exact multiplicities,
support and dominance, with the weight-preserving splitting specified
below. For a positive composition $\kappa=(\kappa_1,\ldots,\kappa_m)$,
the multiplicity is
\[
N(\kappa)=\prod_{i=1}^{m-1}\binom{\kappa_i+\kappa_{i+1}-1}{\kappa_{i+1}}.
\]
A Dyck-path bijection gives a self-contained proof. We characterize all
nonzero matrix entries, recover their compositions explicitly, and prove
prefix dominance, equal hike and upper triangularity. The printed
odd-length splitting omits its final odd part; adding that part restores
weight and reproduces the source's Figure~8. The literal fixed-weight
matrix also satisfies the conjectured inequalities, with a smaller support.
Further corollaries give explicit moment coefficients, canonical lifts,
positive diagonal entries and support counts at each hike.
\end{abstract}
\noindent\textbf{License.} Creative Commons Attribution 4.0 International
(CC BY 4.0).

\section{The composition multiplicity}
The posing source~\cite[\S8.5, pp.~46--47]{PS} conjectures upper
triangularity, prefix dominance and equal hike for its multiplicity matrix
(Conjecture~8.22), and asks for composition multiplicities, individual
matrix entries and exact vanishing conditions (Problem~8.23).
The printed statement begins ``The matrix $N_n$ is upper triangular.''
Theorem~\ref{thm:product} identifies its classical enumerative ingredient;
Theorem~\ref{thm:support} supplies the source-specific inverse and support.

Write $\Comp(n)$ for the positive compositions of $n$ and $\NC(n)$ for
the noncrossing set partitions of $\{1,\ldots,n\}$. Draw an arc between
successive elements of each block. Immediately before vertex $i$, an
arc $(a,b)$ is active when $a<i\le b$. Let $H_i$ be the number of active
arcs. The statistics in Petrov--Scott~\cite[\S8.1, (8.1)]{PS} are
\[
 \ell_j=\#\{i:H_i=j\},\qquad
 g_j=\#\{i:i\text{ has a predecessor of rank }j
                 \text{ among active left endpoints}\}.
\]
Ranks start at one, in increasing order of the left endpoints.
For a noncrossing partition, a vertex with a predecessor closes or
continues the innermost active block. Its predecessor has rank $H_i$.
With $p=\max_i H_i$, Definition~8.16 of~\cite{PS} therefore gives
\[
 z(\pi)=(\ell_0,g_1,\ell_1-g_1,\ldots,g_p,\ell_p-g_p),
 \tag{1}
\]
with the last coordinate deleted precisely when it is zero. For $p=0$
this means $z(\pi)=(n)$. Define
$N(\kappa)=\#\{\pi\in\NC(n):z(\pi)=\kappa\}$, as in
Corollary~8.20 and Problem~8.23 of~\cite{PS}.

\begin{theorem}\label{thm:product}
For every $\kappa\in\Comp(n)$,
\[
 N(\kappa)=\prod_{i=1}^{m-1}
       \binom{\kappa_i+\kappa_{i+1}-1}{\kappa_{i+1}}.
 \tag{2}
\]
The empty product is one. In particular every composition has positive
multiplicity.
\end{theorem}
\begin{proof}
Process the vertices in order, retaining a stack of open blocks.
A first but not last block element opens a block; a last but not first
element closes the top block; an interior block element extends the top
block; and a singleton starts and finishes its own block.
Noncrossing is exactly what makes the stack rule applicable. These four
possibilities give an up-step, a down-step, a transient horizontal step,
and a singleton horizontal step, respectively. Transients are forbidden
at height zero. Conversely these steps reconstruct a unique partition:
an up-step pushes a new block, a down-step appends its vertex and pops,
a transient appends its vertex to the top block without popping, and a
singleton creates a one-element block. The stack guarantees noncrossing.

Replace the four step types, respectively, by
\[
 UU,\qquad DD,\qquad DU,\qquad UD.
 \tag{3}
\]
This produces a Dyck path of semilength $n$. Every Dyck path has a unique
inverse by grouping consecutive steps in pairs; a $DU$ pair at even
height zero is impossible. Thus (3) is a bijection.

Let $U_j,D_j,S_j,T_j$ count the original up-, down-, singleton and
transient steps starting at height $j$. Then
\[
 \ell_j=U_j+D_j+S_j+T_j,\quad g_j=D_j+T_j,\quad
 U_{j-1}=D_j,
 \tag{4}
\]
with $D_0=T_0=0$. In the doubled path, up-steps ending at height $2j+1$
number $U_j+S_j=\ell_j-g_j$; up-steps ending at height $2j$ number
$U_{j-1}+T_j=g_j$. At height one the count is $U_0+S_0=\ell_0$.
Consequently the numbers of up-steps ending at successive positive
heights are exactly the coordinates of (1).

Under the contour bijection between Dyck paths and rooted ordered
trees, these up-steps create the vertices at successive depths.
There are $\kappa_i$ ordered vertices at depth $i$. Distributing the
$\kappa_{i+1}$ ordered children among them is a weak composition into
$\kappa_i$ slots, giving
$\binom{\kappa_i+\kappa_{i+1}-1}{\kappa_{i+1}}$ choices.
The root has its $\kappa_1$ children in one way, the last generation has
no children, and choices at intermediate depths are independent.
Their product proves (2). This is the tree-profile interpretation of
Flajolet's Stieltjes--Rogers coefficient formula~\cite[p.~132,
Proposition~3B; pp.~150--151, Proposition~10(i)]{F}.
Flajolet's Proposition~10(i), footnote~2, attributes an independent
derivation to the earlier work of R.~C.~Read (1979)~\cite{Read}.
This attribution to Read is made through Flajolet.
The same profile product was stated again in a 2015 MathOverflow
question; Grinberg supplied an excursion-insertion proof, and Gessel
identified Flajolet's Proposition~3B as its source~\cite{MO}.
\end{proof}

\section{Weight-preserving splitting and dominance}
A Fibonacci word is a word over $\{1,2\}$; its weight is the sum of its
letters and its hike $\hik$ is its number of twos. For a word with
$p\ge1$ twos, write
$w=1^{r_0}2\,1^{r_1}2\cdots2\,1^{r_p}$ and put
\[
 \sigma(w)=(r_0+1,r_1+2,\ldots,r_{p-1}+2,r_p+1).
 \tag{5}
\]
For $w=1^n$, put $\sigma(w)=(n)$. These are precisely the positive
compositions with interior coordinates at least two. Their inverse is
denoted $\Fib$. Both maps preserve weight.

For $m\ge2$, set $p=\lfloor m/2\rfloor$ and, when $m=2p$, put
$\kappa_{2p+1}=0$. Define
\begin{align}
 A(\kappa)&=(\kappa_1,\kappa_2+\kappa_3,
             \ldots,\kappa_{2p-2}+\kappa_{2p-1},
             \kappa_{2p}+\kappa_{2p+1}),\notag\\
 B^*(\kappa)&=\left(1-p+\sum_{j=0}^{p}\kappa_{2j+1},
                    1+\kappa_2,\ldots,1+\kappa_{2p-2},\kappa_{2p}\right).
 \tag{6}
\end{align}
For $p=1$ the displayed intermediate lists are empty. For $m=1$ set
$A(\kappa)=B^*(\kappa)=\kappa$.

For $m\le3$ and for even $m$, (6) agrees with Definition~8.12
of~\cite{PS}. For odd $m=2p+1\ge5$, the printed first coordinate of
$B$ sums the first $p$ odd parts rather than all $p+1$ of them.
That definition calls $A,B$ ``two auxiliary Fibonacci compositions of
$n$''; its claimed weight for $B$ fails in these odd lengths.
It follows by addition that
\[
 \sum B(\kappa)=n-\kappa_{2p+1},\qquad
 \sum B^*(\kappa)=n.
 \tag{7}
\]
Thus (6) adds $\kappa_{2p+1}$ to that first coordinate. Both the
14 August 2025 preprint and the journal PDF updated on 5 February 2026
retain the former bound. The associated published
corrigendum~\cite{PScorr} concerns notation rendering and does not
supply another splitting formula. For instance $\kappa=(1,1,1,1,1)$ gives
$A=(1,2,2)$, $B=(1,2,1)$ and $B^*=(2,2,1)$, corresponding to
$u=221$, literal $v=22$ and corrected $v=122$. Formula (2) gives
multiplicity one, matching the $(221,122)$ entry of Figure~8.

For two words of weight $n$ define $N^*_{u,v}=N(\kappa)$ when
$(u,v)=(\Fib(A(\kappa)),\Fib(B^*(\kappa)))$ for a composition
$\kappa\in\Comp(n)$, and zero otherwise. The uniqueness needed for
this definition follows in Theorem~\ref{thm:support}.
Write $u\succeq v$ when the sum of the first $q$ letters of $u$ is at
least the corresponding sum for $v$, for every common prefix length $q$.

\begin{theorem}\label{thm:dominance}
If $N^*_{u,v}\ne0$, then $\hik(u)=\hik(v)$ and $u\succeq v$.
The matrix $N^*_n$ is upper triangular when its rows and columns are
listed in descending lexicographic order.
\end{theorem}
\begin{proof}
Let $a=A(\kappa)$ and $b=B^*(\kappa)$. For $m=1$ the words coincide.
Otherwise $a,b$ have $p+1$ coordinates and weight $n$, so both words
have hike $p$ and length $n-p$. Equation (6) gives
\[
 a_i-b_i=\kappa_{2i-1}-1\ge0\quad(2\le i\le p),\qquad
 a_{p+1}-b_{p+1}=\kappa_{2p+1}\ge0.
 \tag{8}
\]
Since their total sums agree, for $1\le j\le p$ we have
$\sum_{i\le j}a_i\le\sum_{i\le j}b_i$.
By (5) the position of the $j$th two in $\Fib(a)$ is
$\sum_{i\le j}a_i-j+1$, so each two of $u$ occurs no later than the
corresponding two of $v$. Every prefix of $u$ therefore contains at least
as many twos as the corresponding prefix of $v$; adding the prefix
length proves dominance. At the first differing letter, dominance
forces $u$ to have 2 and $v$ to have 1. Thus $u$ precedes $v$ in the
stated ordering, proving upper triangularity.
\end{proof}

\begin{corollary}\label{cor:literal}
The same dominance, hike and triangularity assertions hold for the
literal fixed-weight matrix obtained by using the printed $B$ and
requiring both indices to have weight $n$.
\end{corollary}
\begin{proof}
Equation (7) excludes exactly the odd compositions of length at least
five. For every remaining composition, $B=B^*$, so the literal matrix
has a subset of the support of $N^*_n$ with the same entries.
\end{proof}

\section{Exact support and individual entries}
\begin{theorem}\label{thm:support}
Let $u,v$ have weight $n\ge1$. If their hikes differ, $N^*_{u,v}=0$.
If their common hike is zero, $u=v=1^n$ and $N^*_{u,v}=1$.
For common hike $p\ge1$, write $a=\sigma(u)$ and $b=\sigma(v)$.
Then
\[
 N^*_{u,v}>0\quad\Longleftrightarrow\quad
 a_i\ge b_i\quad(2\le i\le p+1).
 \tag{9}
\]
For a supported pair its unique composition is
\begin{align}
 \kappa_1&=a_1,\notag\\
 \kappa_{2j}&=b_{j+1}-1,\qquad
 \kappa_{2j+1}=a_{j+1}-b_{j+1}+1\quad(1\le j<p),\notag\\
 \kappa_{2p}&=b_{p+1},
 \tag{10}
\end{align}
followed by $\kappa_{2p+1}=a_{p+1}-b_{p+1}$ exactly when this
difference is positive. Its entry is given by (2).
\end{theorem}
\begin{proof}
Necessity is (8). Conversely (9) makes every part of (10) positive:
$b_{j+1}\ge2$ for $j<p$, and the endpoints are positive. The optional
last part is positive by its inclusion rule. Summing the parts gives
$\sum_i a_i=n$. Substitution into (6) recovers all coordinates of
$a$ and all coordinates of $b$ except possibly the first. Equality of
total weights then recovers the first coordinate too. Formula (2)
makes the resulting entry positive. Finally (6) forces every part in
(10) and forces whether the last part is present, proving uniqueness.
\end{proof}

For the literal matrix, add to (9) the condition
$a_{p+1}=b_{p+1}$ when $p\ge2$; when $p=1$ no additional condition
is needed. This is an exact vanishing criterion for both interpretations
of the matrix.

\begin{corollary}
The corrected matrix has $2^{n-1}$ nonzero entries, and
\[
 \sum_{\kappa\in\Comp(n)}N(\kappa)
 =\sum_{u,v}N^*_{u,v}=\frac1{n+1}\binom{2n}{n}.
 \tag{11}
\]
\end{corollary}
\begin{proof}
Theorem~\ref{thm:support} gives an injective correspondence from the
$2^{n-1}$ positive compositions to supported pairs, and
Theorem~\ref{thm:product} gives positive coefficients.
The Dyck paths in its proof are partitioned by their height profiles.
There are $\binom{2n}{n}-\binom{2n}{n-1}=\binom{2n}{n}/(n+1)$
Dyck paths, by reflecting a path up to its first step below zero.
\end{proof}

\section{Sibling consequences and sharper boundaries}
In this section put $\kappa_{2p+1}=0$ when $m=2p$.
For a Fibonacci composition $b$ of length $p+1$, write
$\operatorname{dep}(b)=(b_2-1,\ldots,b_p-1,b_{p+1})$;
for $p=0$ this is the empty list.

\begin{corollary}[Statistics and a canonical lift]\label{cor:lift}
For $z(\pi)=\kappa$ one has
$\ell(\pi)=A(\kappa)$ and
$g(\pi)=(\kappa_2,\kappa_4,\ldots,\kappa_{2p})
=\operatorname{dep}(B^*(\kappa))$.
The map $z:\NC(n)\to\Comp(n)$ is surjective and has an explicit
section: at every depth, give all vertices of the next depth to the
first vertex of the current depth, then apply the contour and paired-step
inverse from the proof of Theorem~\ref{thm:product}.
\end{corollary}
\begin{proof}
The statistics follow by substituting (1) into (6).
For any positive composition, the prescribed rooted ordered tree exists:
each nonfinal depth has at least one vertex. Its depth profile is
$\kappa$, so its inverse partition has $z(\pi)=\kappa$.
This reproves the statistics and surjectivity of~\cite[Lemma~8.17,
Proposition~8.18]{PS} and gives another canonical lift alongside its
Remark~8.19. Depletion ignores the first coordinate, so replacing $B$
by $B^*$ does not change $g$ or the moment monomial.
\end{proof}

\begin{corollary}[Explicit moment coefficients]\label{cor:moments}
Define $M_n=\sum_{\pi\in\NC(n)}
\prod_{j\ge0}c_{j+1}^{\ell_j(\pi)}
\prod_{j\ge1}t_j^{g_j(\pi)}$ for commuting indeterminates $c_j,t_j$.
Then
\[
 M_n=\sum_{\kappa\in\Comp(n)}
 \left(\prod_{i=1}^{m-1}
 \binom{\kappa_i+\kappa_{i+1}-1}{\kappa_{i+1}}\right)
 c_1^{\kappa_1}\prod_{j=1}^{p}
 c_{j+1}^{\kappa_{2j}+\kappa_{2j+1}}t_j^{\kappa_{2j}}.
 \tag{12}
\]
Under the totally positive specialization of~\cite[Proposition~8.8,
Corollary~8.20]{PS}, $M_n$ is its $n$th Stieltjes moment.
\end{corollary}
\begin{proof}
Group partitions by $z(\pi)$, use Corollary~\ref{cor:lift} for the
exponents, and apply (2) to the size of each fiber. This makes the
coefficients of the source's moment expansion explicit; their underlying
Stieltjes--Rogers formula is Flajolet's~\cite[Proposition~3B]{F}.
\end{proof}

\begin{corollary}[Diagonal, hike blocks and support counts]\label{cor:counts}
Every diagonal entry of both matrices is positive. Their rows and columns
are nonempty, and the matrices are invertible over $\mathbb Q$.
They have zero entries between distinct hikes. For $p\ge1$, the corrected
hike-$p$ block has exactly $\binom n{2p}$ nonzero entries, with
$\binom nr=0$ for $r>n$. The hike-zero block has one entry.
For $n\ge2$ the literal matrix has $2^{n-2}+1+\binom{n-1}{2}$
nonzero entries; at $n=1$ it has one.
\end{corollary}
\begin{proof}
In (10), take $a=b=\sigma(w)$. This gives $\kappa=(n)$ if $p=0$,
and otherwise an even-length composition with first part $a_1$,
last part $a_{p+1}$, and intermediate pairs $(a_{j+1}-1,1)$ for
$1\le j<p$. It splits to $(w,w)$ even for the printed $B$.
Thus (2) gives positive diagonal entries, both projections of the
splitting are surjective as in~\cite[Lemma~8.14]{PS}, and triangularity
gives a nonzero determinant. Distinct hikes have zero entries by
Theorem~\ref{thm:support}.

Hike $p\ge1$ corresponds precisely to composition lengths $2p$ and
$2p+1$. Their counts sum to
$\binom{n-1}{2p-1}+\binom{n-1}{2p}=\binom n{2p}$.
For the literal matrix only length one, length three and even lengths
remain. For $n\ge2$, the even lengths contribute $2^{n-2}$ (pair subsets
of the $n-1$ separator positions by toggling one fixed position),
and lengths one and three contribute $1$ and $\binom{n-1}{2}$.
\end{proof}

\begin{corollary}[The two one-alphabet specializations]\label{cor:alphabets}
Setting all $c_j=1$, or setting all $t_j=1$, in (12) leaves exactly
$\#\mathrm{YF}_n$ distinct monomials, each with positive integer
coefficient, as observed in~\cite[\S8.5, opening paragraph]{PS}.
\end{corollary}
\begin{proof}
For $t_j=1$, the exponent list is $A(\kappa)=\sigma(u)$.
For $c_j=1$, it is $\operatorname{dep}(\sigma(v))$, which determines
$\sigma(v)$ at fixed weight by recovering its first coordinate.
Every $u$ and every $v$ occurs by the diagonal construction in
Corollary~\ref{cor:counts}; all coefficients in (12) are positive.
\end{proof}

\begin{corollary}[Dominance does not characterize support]\label{cor:converse}
The converse of Conjecture~8.22 is not asserted in the source and fails:
$u=221$ dominates $v=212$ and both have hike two, but
$N^*_{221,212}=N_{221,212}=0$.
\end{corollary}
\begin{proof}
The prefix sums are $(2,4,5)$ and $(2,3,5)$.
The row compositions are $(1,2,2)$ and $(1,3,1)$, so (9) fails
at their second coordinate. The literal support is smaller.
\end{proof}

\section{Coverage and verification scope}
\begin{center}\small
\begin{tabular}{@{}p{.22\textwidth}p{.7\textwidth}@{}}
\toprule
Source item in~\cite{PS} & Mathematical coverage\\
\midrule
Problem 8.23: $N(\kappa)$ & Answered: yes, by (2).
Reproved; prior coefficient formula Flajolet (1980), Proposition 3B.\\
Problem 8.23: entries and zeros & Answered: yes, by
Theorem~\ref{thm:support}, (10), and the literal variant.\\
Conjecture 8.22 & Proved: Theorem~\ref{thm:dominance} for (6),
Corollary~\ref{cor:literal} for the literal fixed-weight matrix.\\
Definition 8.12: weight claim & Refuted-and-corrected for odd $m\ge5$:
the five-part example and (7) identify the deficit; (6) restores weight.\\
Lemma 8.14 & Reproved; prior proof Petrov--Scott (2025).
Theorem~\ref{thm:support} gives injectivity, Corollary~\ref{cor:counts}
gives both surjections for the corrected map.\\
Lemma 8.17; Proposition 8.18 & Reproved; prior proofs Petrov--Scott
(2025). Corollary~\ref{cor:lift} gives the statistics and surjectivity.\\
Remark 8.19 & Alternative canonical lift in Corollary~\ref{cor:lift};
existence was already established by Petrov--Scott (2025).\\
Corollary 8.20 & Reproved with explicit coefficients in (12);
prior moment expansion Petrov--Scott (2025), coefficient formula
Flajolet (1980).\\
\S8.5, opening observation & Reproved; prior observation Petrov--Scott
(2025). Corollary~\ref{cor:alphabets} counts both specialized alphabets.\\
Figure 8 & Finite check: all 64 entries agree with (2), (6), (10).\\
Converse; refined counts & Not posed; obstruction proved in
Corollary~\ref{cor:converse}; counts and invertibility proved in
Corollary~\ref{cor:counts}.\\
Other questions in \S18 & Not covered: Problems 4.10, 7.4, 7.5, 10.5;
Conjecture 6.3; promotion, truncations, and involution directions.\\
\bottomrule
\end{tabular}
\end{center}

The Lean declaration
\href{https://github.com/the-omega-institute/trureturing/blob/e277c80d045570daf43bd03b4e313f7ce457b4e2/D5/S3/Combinatorics/YoungFibonacci/PetrovScottMultiplicityDominance.lean}{\texttt{PetrovScottMultiplicityDominance.}\allowbreak\texttt{result}}
states the literal fixed-weight conclusion of
Corollary~\ref{cor:literal}, using the printed $B$ and its defined
noncrossing-partition multiplicity. Its hypotheses require both words to
have the same positive weight; it does not state the corrected-matrix
claim or the product formula. The declaration uses
\href{https://github.com/the-omega-institute/trureturing/blob/e277c80d045570daf43bd03b4e313f7ce457b4e2/D5/S1/Words/UnaryGapWords.lean}{\texttt{UnaryGapWords.gap}}
to concatenate runs of ones separated by twos in its reverse-word
comparison. This is a word constructor, not a multiplicity theorem.
The product formula, correction,
corrected-matrix theorem and exact support are established by the
written proofs above. The accompanying exact-integer checker independently
enumerates set partitions and Dyck paths through $n=9$, word pairs
through $n=12$, and compositions through $n=16$. It also checks the
moment exponents, canonical lifts, diagonal entries and refined support
counts, and the explicit converse counterexample. These finite checks
complement, rather than replace, the all-$n$ proofs.

\begin{thebibliography}{9}
\small
\bibitem{PS} L. Petrov and J. Scott,
\emph{Random Fibonacci words via clone Schur functions},
Forum of Mathematics, Sigma \textbf{14} (2026), e15, 1--90.
\href{https://doi.org/10.1017/fms.2025.10165}{doi:10.1017/fms.2025.10165}.
Posing locators: Definitions 8.12, 8.16 and 8.21, Conjecture 8.22,
Problem 8.23; also
\href{https://arxiv.org/abs/2412.21126v2}{arXiv:2412.21126v2}
(14 August 2025), pp. 43--47.
\bibitem{PScorr} L. Petrov and J. Scott,
\emph{Random Fibonacci Words via Clone Schur Functions---Corrigendum},
Forum of Mathematics, Sigma \textbf{14} (2026), e46, 1.
\href{https://doi.org/10.1017/fms.2026.10186}{doi:10.1017/fms.2026.10186}
(27 March 2026).
\bibitem{F} P. Flajolet,
\emph{Combinatorial aspects of continued fractions},
Discrete Mathematics \textbf{32} (1980), 125--161,
especially p. 132, Proposition 3B.
\href{https://doi.org/10.1016/0012-365X(80)90050-3}{doi:10.1016/0012-365X(80)90050-3}.
\bibitem{Read} R. C. Read,
\emph{The chord intersection problem},
Annals of the New York Academy of Sciences \textbf{319} (1979), 444--454.
\href{https://doi.org/10.1111/j.1749-6632.1979.tb32822.x}{doi:10.1111/j.1749-6632.1979.tb32822.x}.
Attribution via Flajolet~\cite[Proposition~10(i), footnote~2]{F}.
\bibitem{MO} D. Grinberg and I. Gessel, answers to
\emph{Number of Dyck paths with prescribed number of edges},
MathOverflow, 20 February 2015.
\href{https://mathoverflow.net/a/198057}{Grinberg's insertion proof};
\href{https://mathoverflow.net/a/198059}{Gessel's reference to Flajolet}.
\end{thebibliography}
\end{document}
